\section{Alzheimer's Disease and Biomarkers}

Alzheimer’s disease (AD) has traditionally been diagnosed based on clinical symptoms, with diagnostic criteria based on progressive decline in memory and other cognitive functions~\cite{McKhann1984}. 
Mild cognitive impairment (MCI) describes a condition in which cognitive decline is present but does not significantly affect independence in daily activities~\cite{Albert2011MCI}. 
MCI can occur before dementia, but it can have different causes and does not always indicate underlying AD pathology~\cite{Albert2011MCI}. 
Therefore, clinical groups such as healthy controls (HC), MCI, and AD describe differences in cognitive and functional state rather than the underlying biological processes of the disease.

The main pathological features of AD are the accumulation of amyloid-beta (A$\beta$) and tau proteins.
A$\beta$ accumulates outside neurons and forms amyloid plaques, while abnormal tau accumulates inside neurons and forms neurofibrillary tangles~\cite{Busche2020Synergy}. 
These protein accumulations are considered key pathological features of AD\@.

A$\beta$ and tau are closely linked to the development and progression of AD\@.
Changes in these biomarkers can begin before clear cognitive symptoms appear.
A$\beta$-related changes can occur relatively early in the disease process, followed by further changes related to neuronal injury and cognitive decline.
A$\beta$ and tau are also not fully separate processes, as evidence suggests that they interact and together contribute to neuronal dysfunction~\cite{Busche2020Synergy}.
Their effects on brain activity may also differ across clinical stages.
Whole-brain modelling has shown a stronger influence of A$\beta$ on brain dynamics in MCI, while tau had a stronger influence in AD~\cite{Patow2023}.
The combination of both A$\beta$ and tau also provided a better representation of empirical brain dynamics than either biomarker individually~\cite{Patow2023}.

The use of biomarkers has led to an important distinction between the clinical diagnosis and biological definition of AD\@.
The NIA-AA Research Framework introduced the AT(N) classification system, where (A) represents A$\beta$ pathology, (T) represents tau pathology, and ((N)) represents neurodegeneration or neuronal injury~\cite{Jack2018}.
Under this framework, AD can be defined based on its underlying biological pathology rather than clinical symptoms alone~\cite{Jack2018}.
Clinical groups such as HC, MCI, and AD and biological measures of A$\beta$ and tau therefore provide different but related information.
An individual's clinical state does not always directly correspond to their biological AD status~\cite{Jack2018}.